\chapter{Conclusion} 
\label{cha:conclusion}

The main insights of this investigation are...
\begin{infobox}[][Conclusion vs. summary]
In the final chapter of the manuscript, you should provide a conclusion of the work. However, in practice many authors formulate a summary instead of a genuine conclusion. Hence, please note the differences:
\begin{itemize}
    \item \textbf{Conclusion:} Is meant to provide a final interpretation of the study's results and offer insights into the broader implications of the findings. It often discusses the significance of the research in the context of existing knowledge. If relevant, key quantative figures should be included. In simple words: what are the take-home insights of the work?
    \item \textbf{Summary:} Aims to consecutively recap all key points of the entire work. It covers a step-by-step overview of the main objectives, utilized assumptions, applied methods, empirical experiments and the found results. The summary therefore goes beyond the conclusion, as it focus not the findings, but provides a larger overview of the entire work.
\end{itemize} 
In this template, the manuscript already starts with an abstract which is a very concise and brief summary. Therefore, the final conclusion should not(!) be a second, redundant summary.   
\end{infobox}

% Examples: correct vs. incorrect conclusion
\begin{successbox}[][Correct conclusion -- example]
This thesis shows that the proposed approach improves robustness under realistic noise while meeting real-time constraints. The key takeaway is that combining temporal context with lightweight uncertainty calibration yields consistent gains up to \SI{5.2}{\percent} across diverse datasets. This implies practical deployability on embedded hardware with reduced maintenance effort. Remaining limitations include sensitivity to extreme domain shifts by \SI{42}{\percent}, motivating follow-up work on adaptive calibration.
\end{successbox}

\begin{errorbox}[][Incorrect conclusion (summary) -- example]
In Chapter~1 we introduced the problem, in Chapter~2 we reviewed related work, in Chapter~3 we described the method, in Chapter~4 we presented experiments, and in Chapter~5 we discussed the results. The dataset and training details were also explained. This restates the thesis structure but does not synthesize what the results mean or why they matter; it is a summary, not a conclusion.
\end{errorbox}

Another hint: to ensure that the conclusion is concise and to the point, one can consider using a bullet point list and avoid long, complex sentences. This will help the reader to extract the main insights quickly.


\section{Future work}
\label{sec:futurework}
The outlook of this work should provide a brief overview of potential future research directions. This can include open questions, unsolved problems, or potential improvements of the current work. It is also possible to discuss potential applications of the findings in practice. A bullet point list is again a good choice for this section to ensure a clear and concise presentation of the future work. 

\section{Publish your work}
\label{sec:publish}
We encourage students to voluntarily publish their work through the University Siegen library. Visit the following link for more information on the publishing process: \href{https://www.ub.uni-siegen.de/en/search-tools-and-catalogues/eresources/opus-siegen/information-for-authors/publication-types/bachelor-and-master-theses/}{https://www.ub.uni-siegen.de/en/search-tools-and-catalogues/eresources/opus-siegen/information-for-authors/publication-types/bachelor-and-master-theses/}.

If you like to publish your work, follow this official process first. Refrain from publishing your work on other third-party platforms (e.g., ResearchGate, Academia.edu, etc.) before the finalizing the official publication process! 